\paragraph{明确 SFibAI 基线身份。}
方法来源为 \citet{xu2026sfibai}。基线采用作者的公开实现。基线使用 ResNet-50、全局平均池化、线性 $2048\rightarrow36$ 头、softmax 与 0.0--3.5 上的期望分数，没有切面和定位预测头。上游允许指定初始化检查点，未提供时使用随机初始化骨干；共同 ImageNet 初始化是针对本次比较的适配，不是上游默认值。

\paragraph{评分目标。}
VCRE-Fib 与 SFibAI 采用相同的纤维化评分尺度和混合损失权重 $(\alpha,\beta,\gamma)=(1,0.02,0.02)$。对于目标 bin $j\in\{0,\ldots,35\}$ 和预测概率 $p_k$，两者均先将 bin 索引转换为真实临床分值，再计算 MSE 与临床等级边界惩罚：
\[
\begin{aligned}
y&=0.1j,\qquad \hat y=0.1\sum_{k=0}^{35}kp_k,\\
L_{\mathrm{grade}}&=L_{\mathrm{KL,local}}+0.02(\hat y-y)^2
 +0.02|\hat y-y|\mathbf{1}(\hat c\ne c).
\end{aligned}
\]
其中 $c$、$\hat c$ 分别由 $y$、$\hat y$ 按 0.5、1.5、2.5 阈值映射为四级，分数恰好等于阈值时归入较高等级。上式为逐图损失，各项在 batch 内取均值。MSE 和绝对误差边界项均在 0.0--3.5 的真实分值尺度上计算。

局部 KL 采用定义在全部 36 个 bin 上、标准差为一个 bin 的归一化高斯软目标，并在以目标 bin 为中心的局部窗口内计算。两方法采用这一混合评分损失形式，评分单位及各分项权重保持一致。

\paragraph{训练协议。}
四模型均采用相同随机种子（2026）、ImageNet-1K-V2 权重、batch size 24、bfloat16、AdamW 学习率与权重衰减均为 $10^{-4}$，以及每 15 epoch 将学习率乘以 0.6 的 StepLR。ROI 预处理、拉伸缩放、直角旋转/饱和度增强、患者互斥划分及上述评分目标相同。这统一了公开 SFibAI 中不同的裁剪、增强和初始化默认值。SFibAI、两个删除变体和完整 VCRE-Fib 的训练轨迹及检查点比较均使用相同的 90 epoch 上限。在验证集上从第 21--90 epoch 中依次按较低 $\risk$、较低 image COR、较早 epoch 的顺序选择检查点，替换上游容差准确率选择器。四者所选 epoch 依次为 65、80、39、47。本比较使用重新整理的队列，不声称复现 SFibAI 原始队列实验。

\paragraph{VCRE-Fib 的新增监督。}
六分类切面交叉熵权重为 0.1。弱框损失结合框内 log-sum-exp 池化（温度 10）和框外响应占比，内部权重为 1，总权重为 0.1。仅有效、非空且目标 bin $j>0$ 的弱框参与；排除精确分数 0.0，不排除整个 F0。弱目标在 ROI 坐标系中变换，支持图通过评分学习。辅助损失更新共享编码器，仅在评分接口截断 $q$ 与 $a$。w/o View 删除切面头、六个适配器、后验混合和切面 CE，保留无切面条件的区域证据与弱定位。w/o Weak Loc.\ 删除定位头和弱框损失，使用中性 $a_u=1$ 并保留切面与支持路径。两个删除变体均由规定初始化独立训练。

